\documentclass[../../../include/open-logic-section]{subfiles}

\begin{document}

\olfileid{sth}{ord-arithmetic}{mult}

\olsection{الضرب الترتيبي}

نأتي الآن إلى الضرب الترتيبي، ونسلك فيه قريبًا من مسلك الجمع. فلو
أردنا ضرب~$\alpha$ في~$\beta$، تصورنا شبكة مستطيلة عرضها~$\alpha$
وارتفاعها~$\beta$. وحاصل ضرب~$\alpha$ و$\beta$ هو ما ينتج من اجتياز صف كامل، ثم الصف
الذي يليه، وهلم جرًا حتى تستوفى الشبكة. وبعبارة جامعة، ينتج حاصل ضرب
$\alpha$ و$\beta$ بإحلال نسخة من~$\alpha$ محل \emph{كل} عنصر في
$\beta$.

ولصياغة المعنى صوريًا نستعمل الترتيب المعجمي المعكوس على حاصل الضرب
الديكارتي لـ$\alpha$ و$\beta$.

\begin{defn}
لأي عددين ترتيبيين~$\alpha, \beta$، يكون حاصل ضربهما
$\alpha \ordtimes \beta = \ordtype{\alpha \times \beta, \rlexless}$.
\end{defn}
\noindent
ولا بد، مرة أخرى، من إثبات سلامة التعريف.

\begin{lem}\ollabel{ordtimeslessiswo}
$\tuple{\alpha \times \beta, \rlexless}$ ترتيب جيد، لأي عددين
ترتيبيين~$\alpha$ و$\beta$.
\end{lem}

\begin{proof}
الحجة عين ما في~\olref[add]{ordsumlessiswo}.
\end{proof}
\noindent
ثم توصف حركة الضرب بالمعادلات الآتية.

\begin{lem}\ollabel{ordtimesrecursion}
لأي عددين ترتيبيين~$\alpha, \beta$:
\begin{align*}
	\alpha \ordtimes 0 &= 0\\
	\alpha \ordtimes (\beta \ordplus 1) &=
		(\alpha \ordtimes \beta) \ordplus \alpha\\
	\alpha  \ordtimes \beta &=
		\supstrict_{\delta < \beta}(\alpha \ordtimes \delta) &&
		\text{عندما يكون $\beta$ عددًا ترتيبيًا حديًا}.
\end{align*}
\end{lem}

\begin{proof}
متروك تمرينًا.
\end{proof}

وكما جاز في الجمع، كان يجوز أن نتخذ هذه المعادلات العودية نفسها
تعريفًا للضرب الترتيبي بدل التعريف المباشر. وكذلك يشترك الضرب مع
المألوف في الخواص الآتية.

\begin{lem}\ollabel{ordinalmultiplicationisnice}
إذا كانت~$\alpha, \beta, \gamma$ أعدادًا ترتيبية، فإن:
\begin{enumerate}
	\item\ollabel{ordtimes1} إذا كان~$\alpha \neq 0$ و$\beta < \gamma$،
	فإن~$\alpha \ordtimes \beta < \alpha \ordtimes \gamma$؛
	\item\ollabel{ordtimes2} إذا كان~$\alpha \neq 0$ و
	$\alpha \ordtimes \beta = \alpha\ordtimes\gamma$، فإن
	$\beta = \gamma$؛
	\item\ollabel{ordtimes3} $\alpha \ordtimes (\beta \ordtimes
	\gamma) = (\alpha \ordtimes \beta) \ordtimes \gamma$؛
	\item\ollabel{ordtimes4} إذا كان~$\alpha \leq \beta$، فإن
	$\alpha \ordtimes \gamma \leq \beta \ordtimes \gamma$؛
	\item\ollabel{ordtimes5} $\alpha \ordtimes (\beta \ordplus
	\gamma) = (\alpha \ordtimes \beta )\ordplus (\alpha\ordtimes
	\gamma)$.
\end{enumerate}
\end{lem}

\begin{proof}
متروك تمرينًا.
\end{proof}

ولك أن تستخرج غيرها من النتائج ما شئت. على أن
\olref[ord-arithmetic][add]{ordsumnotcommute} يجعل القضية الآتية
متوقعة.

\begin{prop}
الضرب الترتيبي ليس تبادليًا:
$2 \ordtimes \omega = \omega < \omega \ordtimes 2$
\end{prop}

\begin{proof}
$2 \ordtimes \omega = \supstrict_{n < \omega}(2\ordtimes n) = \omega
\in \supstrict_{n < \omega}(\omega \ordplus n) = \omega \ordplus
\omega = \omega \ordtimes 2$.
\end{proof}
\noindent
ووجه ذلك في الحدس بيّن. فحساب~$2 \ordtimes \omega$ يستبدل بكل عدد
طبيعي شيئين؛ وهذا لا يغير نمط الترتيب لو أدخلنا بين كل عددين طبيعيين
عددًا «نصفيًا»، أي نظرنا في الترتيب الطبيعي على
$\Setabs{\nicefrac{n}{2}}{n \in \omega}$. فيظهر حينئذ أن نمط الترتيب
هو نمط~$\omega$ نفسه. أما في~$\omega \ordtimes 2$ فنضع نسختين من
$\omega$ متعاقبتين.

\begin{prob}
أثبت
\olref[sth][ord-arithmetic][mult]{ordtimeslessiswo}،
و\olref[sth][ord-arithmetic][mult]{ordtimesrecursion}،
و
\olref[sth][ord-arithmetic][mult]{ordinalmultiplicationisnice}
\end{prob}

\end{document}
